\chapter{Conclusion \& Future Works}
\label{ch:conclusion}

\section{Summary of Research Findings}
\label{sec:conclusion_summary}
In this capstone research project, we conducted an exhaustive, empirical investigation into log-file anomaly detection across large-scale distributed computing systems and operational telemetry streams. Beginning with an empirical audit of the recent IEEE ICSRS 2025 baseline paper by F{\"a}lt et al.~\cite{falt2025lightweight}, we successfully reproduced their standard PCA baseline ($94.60\%$ vs $94.64\%$ $F_1$) on the Hadoop Distributed File System (\textbf{HDFS\_v1}) by discovering an undocumented $L_2$-normalization preprocessing step on raw event count vectors. Furthermore, by introducing Extreme Value Theory Peaks-Over-Threshold (\textbf{POT}) thresholding, we eliminated the literature's test-set oracle percentile heuristic, achieving $95.95\% \pm 1.47\%$ $F_1$ while adhering strictly to zero test-label leakage. Simultaneously, our 30-seed investigation mathematically demonstrated that their published Robust PCA claim of $90.55\%$ $F_1$ is structurally unreachable (peaking at $87.38\%$) due to low-rank matrix rank-collapse on clean training partitions.

Extending linear models to the massive, multi-node Blue Gene/L (\textbf{BGL}) supercomputing dataset exposed a strict linear ceiling ($\sim 65\%$ $F_1$) caused by high intrinsic matrix rank and asynchronous interleaving across 131,072 processing cores. To solve this, we engineered an end-to-end streaming architecture combining node-centric 6-hour temporal windowing, unsupervised Skip-Gram Word2Vec embeddings, concatenated Mean+Max dual pooling ($z \in \mathbb{R}^{64}$), a deep Multi-Layer Perceptron Autoencoder ($64 \rightarrow 32 \rightarrow 16 \rightarrow 32 \rightarrow 64$), and label-free EVT/POT decision thresholding.

Under our strictly leak-free 10-fold cross-validation protocol reflecting natural operating prevalence ($10.07\%$ anomalies, with every anomaly tested exactly once), our system achieves:
\begin{itemize}
    \item \textbf{Mode A (Headline Label-Free POT):} Precision $93.76\% \pm 2.54\%$, Recall $82.61\% \pm 1.96\%$, and $F_1$-score $\mathbf{87.80\% \pm 1.31\%}$ ($0/10$ fallbacks used).
    \item \textbf{Mode B (Semi-Supervised Validation-Tuned):} Precision $93.06\% \pm 2.59\%$, Recall $84.71\% \pm 1.19\%$, and $F_1$-score $\mathbf{88.66\% \pm 1.07\%}$.
    \item \textbf{Threshold-Free Ranking Quality:} PR-AUC of $\mathbf{91.05\% \pm 0.32\%}$ and ROC-AUC of $\mathbf{97.34\% \pm 0.15\%}$.
\end{itemize}

By evaluating our exact trained models on an artificially balanced test set ($50\%$ anomalies), our models achieved $90.16\%$ (Mode A) and $91.36\%$ (Mode B) $F_1$, conclusively proving that earlier project claims of $\sim 90.44\%$ were inflated by base-rate prevalence and test contamination. Architectural ablations confirmed that the neural autoencoder beats linear PCA by $+27.83\%$ in Mode A $F_1$ and $+6.79\%$ in PR-AUC, while Isolation Forest collapses completely under POT thresholding ($2.07\%$ $F_1$). Finally, hardware benchmarks on consumer Apple M2 CPU silicon proved that the 5,328-parameter model requires only $70.25\text{ KB}$ of memory, executes with $43.2\text{ }\mu\text{s}$ single-window latency, and delivers a pipeline throughput of \textbf{62,076 windows/second} without GPU acceleration.

\section{The Dataset-Dependency Rule}
\label{sec:conclusion_rule}
Synthesizing our empirical results across HDFS and BGL leads to the formulation of the \textbf{Dataset-Dependency Rule} for telemetry anomaly detection:
\begin{enumerate}
    \item \textbf{Structured, Single-Lifecycle Telemetry (e.g., HDFS):} When telemetry is explicitly partitioned by clean transaction or session keys (such as \texttt{block\_id}), normal operational sequences lie on an extremely low-dimensional linear hyperplane (intrinsic rank $\le 2$). For such systems, \textbf{lightweight linear models} (PCA on $L_2$-normalized count vectors with POT thresholding) achieve the absolute empirical ceiling ($\sim 96\%$ $F_1$). Deploying heavy deep sequence networks (such as DeepLog or transformers) is computationally wasteful, financially costly, and offers no performance advantage.
    \item \textbf{Complex, Multi-Tenant Interleaved Telemetry (e.g., BGL):} When telemetry is generated by thousands of concurrent asynchronous nodes without global session keys, the intrinsic subspace rank explodes (rank $239$ on BGL). Linear count models collapse ($\sim 65\%$ $F_1$). To detect anomalies in such complex environments, \textbf{non-linear semantic architectures} (spatial node-centric slicing, continuous dense word embeddings, and deep neural autoencoders) are strictly required to model multi-dimensional coordinate correlations.
\end{enumerate}

\section{Candid Engineering and Operational Limitations}
\label{sec:conclusion_limitations}
In adherence to the highest standards of scientific honesty, we identify four critical limitations of our developed methodology:
\begin{enumerate}
    \item \textbf{Vulnerability to Out-of-Vocabulary (OOV) Concept Drift:} As demonstrated in our chronological evaluation (Section~\ref{sec:results_chronological}), where $68.91\%$ of test tokens were unseen during training, static Word2Vec vocabularies fail when software systems evolve. Skipping OOV tokens causes representations to collapse toward zero, severely degrading temporal generalization.
    \item \textbf{Coarse 6-Hour Window Granularity:} Tumbling 6-hour windows aggregate thousands of individual log entries. While effective for capturing persistent hardware failures, localized transient fault bursts occurring over seconds can be diluted by long intervals of healthy messages within the same window.
    \item \textbf{Static POT Threshold Calibration:} In production environments subject to seasonal or progressive operational load shifts, static POT thresholds calculated from historical validation data lose statistical calibration over time.
    \item \textbf{Upfront Log Parsing Dependency:} Our pipeline throughput benchmark ($62,076\text{ windows/sec}$) measures the latency from pre-extracted template IDs. In live production streaming, the initial Drain log parsing stage introduces additional parsing latency that must be accounted for in end-to-end ingest pipelines.
\end{enumerate}

\section{Ethical Considerations, Carbon Footprint, and Reproducibility}
\label{sec:conclusion_ethics}
The rapid proliferation of massive Large Language Models (LLMs) and transformer architectures in systems engineering has introduced severe environmental and financial overheads, requiring megawatts of power and massive GPU server farms. In contrast, our engineered architecture demonstrates that principled statistical machine learning—combining compact neural autoencoders with Extreme Value Theory—can achieve production-grade detection accuracy while consuming only \textbf{milli-watts of power} on commodity CPU hardware. This represents a significant contribution toward sustainable, low-carbon AI for IT operations (AIOps).

Furthermore, this capstone project strictly adheres to the principles of open science. All evaluation protocols, random seeds, split assertions, and raw results are fully documented and preserved in the project repository, ensuring complete reproducibility without artificial data cherry-picking.

\section{Strategic Directions for Future Work}
\label{sec:conclusion_future_work}
Building upon the empirical findings and limitations identified in this research, several high-impact directions for future inquiry are proposed:
\begin{enumerate}
    \item \textbf{Subword and Byte-Pair Tokenization:} To eliminate the lexical OOV breakdown observed in chronological splits, future architectures should replace discrete template ID embeddings with character-level or subword tokenizers (e.g., FastText or Byte-Pair Encoding) capable of generating continuous representations for previously unseen error messages.
    \item \textbf{Incremental Online Learning:} Implement lightweight, streaming updates to the Word2Vec embedding layer and autoencoder weights, allowing the model to adapt continuously to newly introduced log templates in production without full retraining.
    \item \textbf{Dynamic Sliding-Window EVT Estimation:} Replace static validation POT thresholding with an online, sliding-window EVT estimator that continuously updates the GPD parameters $\xi$ and $\sigma$ over streaming normal error tails, ensuring real-time threshold calibration.
    \item \textbf{Edge and Host Sidecar Daemon Deployment:} Port the serialized $70.25\text{ KB}$ autoencoder runtime into an ultra-low-overhead in-process sidecar daemon written in C++ or Rust, evaluating streaming inference on host operating system daemons and SmartNIC embedded processors.
\end{enumerate}

